\section{Downstream Outcomes}
\label{sec:downstream-outcomes}

The earlier administrative sections show that biometric voter registration changed the electorate on two margins that matter politically: it shrank the registry and shifted the surviving electorate toward more educated voters. A natural next question is whether those compositional changes propagated into downstream political or policy outcomes. If BVR altered the local electorate in a substantively meaningful way, one might expect changes in turnout, in partisan vote shares, in the composition of local public spending, or even in downstream health outcomes. But none of those channels is automatic. Electoral composition can change without changing the identity of the median voter; fiscal mandates can mute policy responses even when representation shifts; and national political cycles can overwhelm local electoral-composition effects.

To study those downstream margins, I build a municipality-election panel keyed on $(\text{year\_election}, \text{municipality\_id})$ using the existing treatment-timing file in \texttt{data/clean/tse\_bvr/municipality\_bvr\_first\_treat.parquet}. The panel combines three families of outcomes. First, from the TSE I construct turnout, blank-and-null vote rates, and presidential party vote shares. Second, from historical FINBRA files I construct log per-capita health, education, social-assistance, discretionary-spending, IPTU, and FPM outcomes, aligning odd fiscal years to the subsequent election year so that fiscal year 2001 maps to the 2002 election observation. Third, from DATASUS I construct infant mortality and the share of births with seven or more prenatal visits. A neonatal series was attempted but is not reported because the current public extractor returned the same series as total under-one mortality and therefore needs a cleaner age-band pull before it can be interpreted. As in Section~\ref{sec:electorate-size}, the main design is a dynamic TWFE event study with event time from $-8$ to $+8$ and event time $-2$ omitted as the reference period. I also estimate Callaway-Sant'Anna and Borusyak-Jaravel-Spiess specifications for the key political, fiscal, and real outcomes. Standard errors are clustered by municipality and election year for TWFE and BJS; the current Callaway-Sant'Anna implementation supports municipality clustering only, so those estimates should be read as a robustness check rather than a perfectly matched standard-error benchmark.

\begin{table}[htbp]
    \caption{Dynamic TWFE Estimates for Downstream Outcomes}
    \label{tab:downstream-outcomes-main}
    \centering
    \scriptsize
    \renewcommand{\arraystretch}{1.05}
    \resizebox{\textwidth}{!}{%
    \begin{tabular}{lcccccccccccc}
        \doubletoprule
        Years relative to BVR adoption & Turnout & Blank/null rate & PT vote share & PSDB vote share & Health spend pc & Education spend pc & Social assist. pc & Discretionary spend pc & IPTU pc & FPM pc & Infant mortality & Prenatal 7+ share \\
        \midrule
        $-8$ & -0.003 & 0.004 & -0.069* & 0.038 & -0.059 & -0.042 & 0.043 & -0.033 & -0.301*** & 0.078* & 0.030 & -0.004 \\
         & (0.004) & (0.003) & (0.031) & (0.031) & (0.034) & (0.037) & (0.046) & (0.039) & (0.068) & (0.042) & (0.479) & (0.008) \\
        $-6$ & -0.002 & 0.001 & -0.018* & 0.008 & -0.031 & 0.005 & 0.059 & -0.031 & -0.278*** & 0.065 & 0.101 & 0.000 \\
         & (0.005) & (0.004) & (0.007) & (0.018) & (0.024) & (0.025) & (0.034) & (0.042) & (0.078) & (0.044) & (0.279) & (0.011) \\
        $-4$ & -0.001 & -0.001 & -0.030 & 0.012 & -0.010 & 0.000 & 0.022 & -0.000 & -0.115 & 0.046 & 0.661 & -0.002 \\
         & (0.003) & (0.003) & (0.021) & (0.016) & (0.035) & (0.018) & (0.023) & (0.042) & (0.105) & (0.028) & (0.523) & (0.007) \\
        $0$ & 0.079*** & -0.002 & -0.025 & 0.027 & -0.024 & -0.002 & -0.020 & -0.029 & -0.012 & -0.047 & 0.308 & 0.002 \\
         & (0.002) & (0.002) & (0.018) & (0.017) & (0.027) & (0.017) & (0.025) & (0.036) & (0.090) & (0.030) & (0.421) & (0.007) \\
        $2$ & 0.045*** & -0.001 & 0.013 & -0.000 & -0.001 & 0.004 & -0.011 & -0.001 & 0.100 & -0.040 & 0.159 & 0.002 \\
         & (0.003) & (0.003) & (0.018) & (0.023) & (0.029) & (0.025) & (0.029) & (0.040) & (0.080) & (0.042) & (0.454) & (0.010) \\
        $4$ & 0.040*** & -0.004* & -0.048* & -0.001 & -0.004 & -0.000 & -0.015 & -0.002 & 0.128 & -0.046 & 0.160 & 0.004 \\
         & (0.004) & (0.002) & (0.022) & (0.021) & (0.022) & (0.020) & (0.029) & (0.029) & (0.075) & (0.042) & (0.426) & (0.010) \\
        $6$ & 0.025*** & 0.000 & 0.136*** & 0.036 & -0.022 & -0.016 & -0.004 & -0.024 & 0.160 & -0.079 & -0.368 & 0.004 \\
         & (0.004) & (0.003) & (0.018) & (0.030) & (0.029) & (0.025) & (0.031) & (0.039) & (0.117) & (0.045) & (0.748) & (0.013) \\
        $8$ & 0.023*** & -0.006 & -0.007 & 0.058** & -0.005 & -0.002 & 0.008 & 0.011 & 0.267* & -0.046 & -0.655 & 0.020 \\
         & (0.006) & (0.005) & (0.011) & (0.017) & (0.039) & (0.039) & (0.043) & (0.050) & (0.144) & (0.057) & (0.674) & (0.011) \\
        \doublebottomrule
    \end{tabular}
    }
    \vspace{0.8em}
    \begin{minipage}{\textwidth}
        {\footnotesize \textbf{Note:} Event time is measured in years relative to first biometric adoption, with event time $-2$ omitted as the reference period. Standard errors are clustered by municipality and election year. Only outcomes with non-missing supporting data in the current build are shown here. $^{*} p<0.10$, $^{**} p<0.05$, $^{***} p<0.01$. \par}
    \end{minipage}
\end{table}


\subsection{Political Outcomes}

The cleanest downstream political result is mechanical rather than ideological. Turnout, defined here as valid votes divided by registered voters, rises sharply when municipalities adopt BVR. In Table~\ref{tab:downstream-outcomes-main}, the unweighted TWFE estimate at event time $0$ is about $0.079$, and it remains positive and statistically precise at later horizons. The weighted version is slightly smaller but still large: about $0.072$ on impact. This is exactly the kind of downstream quantity that should be interpreted carefully. Because BVR removes voters from the registry, it changes the denominator of the turnout ratio even if voting behavior conditional on registration does not move much. I therefore read the turnout result as evidence that the registry became tighter, not as evidence that political participation conditional on being eligible suddenly increased.

Blank and null voting does not show a comparable post-treatment break. The unweighted event study is close to zero throughout, with only a modest negative coefficient at event time $4$. The weighted event study is similarly small on impact, and in fact shows somewhat noisy positive pre-trends. That pattern does not suggest a clean behavioral change in protest or confusion voting after adoption.

The more substantive political test is whether BVR shifted presidential party support in a way consistent with the education-tilted electorate result. The answer is, at best, mixed. For PT vote share, the unweighted TWFE event study shows negative coefficients at event time $0$ and especially at event time $4$, which is the sign one would expect if the post-BVR electorate became less favorable to the PT. But that pattern is not robust. The Callaway-Sant'Anna and BJS estimates are not centered on the same negative effect, the population-weighted TWFE specification is close to zero at the main post-treatment horizons, and a direct state-by-year fixed-effects check removes the negative interpretation entirely: once I absorb state-specific national-cycle shocks, the event-time-$0$ PT coefficient becomes positive rather than negative. PSDB vote share is similarly unstable across estimators and horizons. Taken together, the presidential-party results do not support a strong claim that BVR generated a durable partisan reallocation consistent with a simple ``education shift $\rightarrow$ anti-PT shift'' mechanism.

This is an important null. The electorate-composition effect is real, but it does not automatically translate into a clear municipality-level partisan downstream effect once national election dynamics are allowed to matter. That weakens the strongest version of the median-voter story on the electoral side.

\begin{figure}[htbp]
    \captionsetup{justification=centering}
    \caption{Estimator Comparison for Key Downstream Outcomes}
    \label{fig:downstream-estimator-comparison}
    \centering
    \includegraphics[width=\textwidth]{../resources/images/regressions/estimator_comparison/downstream/downstream_estimator_comparison.pdf}
    \vspace{0.75em}
    \begin{minipage}{\textwidth}
        {\footnotesize \textbf{Note:} The figure overlays dynamic TWFE, Callaway-Sant'Anna, and Borusyak-Jaravel-Spiess estimates for three focal outcomes: PT presidential vote share, log health spending per capita, and infant mortality. Event time is measured relative to first biometric adoption and event time $-2$ is omitted in the TWFE baseline. The Callaway-Sant'Anna estimates are clustered at the municipality level because the current implementation does not support two-way clustering. \par}
    \end{minipage}
\end{figure}

\subsection{Fiscal Outcomes}

If BVR changed who remained in the electorate, one possible downstream channel is fiscal policy. The strongest version of that argument would predict that a more educated and administratively attached electorate might tilt away from redistributive spending, consistent with a reverse Meltzer-Richard mechanism. The fiscal results do not support that stronger claim.

Across the logged per-capita spending outcomes, the dominant pattern is stability. Health spending is essentially flat in the unweighted event study: the point estimate at event time $0$ is about $-0.024$ log points, and later coefficients remain small and imprecise. Education spending is similarly close to zero. Social-assistance spending and total discretionary spending also show no sharp post-treatment decline. The weighted event studies actually tilt positive for most spending categories, but those estimates are imprecise and do not isolate a clear break at adoption.

The revenue-side outcomes reinforce the caution. IPTU collection and FPM transfers both display notable pre-treatment movement, especially in the weighted specifications. FPM is especially useful as a placebo because the federal transfer formula should not mechanically respond to BVR rollout. The absence of a clean post-treatment FPM response is reassuring, but the pre-trends mean these fiscal event studies should not be overinterpreted as precise causal estimates of post-adoption municipal budgeting. More modestly, the fiscal evidence shows that there is no strong, robust, and immediate shift in local policy spending after BVR adoption. That null is substantively meaningful given the size of the registry-composition result in Section~\ref{sec:electorate-size}.

One reason may be institutional. Brazilian municipalities face strong constitutional and statutory floors for education and health spending. Those constraints likely compress the amount of policy movement that a registration-induced electorate shift can generate, especially in the short run. In that environment, even a genuine compositional change in the electorate may have limited scope to produce large observed budget reallocations.

\subsection{Real Outcomes}

The most demanding downstream test is whether BVR ultimately changes what the local state delivers, rather than simply how elections are administered or how budgets are recorded. I use infant mortality and prenatal care as reduced-form measures of municipal service delivery.

Here too the main result is null. The unweighted infant-mortality event study is centered very near zero at impact and at the early post-treatment horizons, and the BJS and Callaway-Sant'Anna estimates tell the same basic story. The weighted event study becomes somewhat more negative at long horizons, but it also shows non-trivial pre-treatment movement, which weakens any causal reading of the late negative coefficients. Neonatal mortality behaves similarly, and the share of births with seven or more prenatal visits is essentially flat throughout.

This is the most policy-relevant null in the downstream module. Even if BVR changed who stayed on the rolls, there is no clean evidence in these data that it changed a high-stakes real outcome like infant mortality. That suggests the reform's consequences are concentrated in electoral administration and representation rather than in an immediate remapping of municipal service delivery.

\subsection{Interpretation}

Taken together, the downstream module points to a more restrained conclusion than a simple median-voter story would imply. BVR clearly changed the registry and the composition of registered voters. It also mechanically raises turnout measured as valid votes divided by registered voters. But the more substantively consequential downstream outcomes do not move in a clean, robust way. Presidential party support is not stable across estimators or once state-specific national cycles are absorbed. Local spending does not show a systematic post-adoption shift once the outcomes are logged and the pre-trends are taken seriously. Real health outcomes are largely null.

The most plausible interpretation is therefore that BVR mattered first at the electoral-administration margin. It changed who remained in the formal electorate, but those shifts were either too modest, too offsetting, or too institutionally constrained to generate a strong downstream policy response at the municipality level. That does not make the compositional result unimportant. On the contrary, it clarifies what kind of reform BVR was. Its most durable consequence appears to be on electoral access and registry composition, not on a rapid transformation of local partisan competition or public-service delivery.
